\chapter{Kết luận và hướng phát triển}

\section{Kết quả đạt được}

Sau quá trình thực hiện, nhóm đã xây dựng và triển khai được hệ thống DocVault với hai phần chính: phần ứng dụng quản lý tài liệu bảo mật và phần DevSecOps cho kiểm thử, triển khai và thu thập bằng chứng vận hành.

Ở phần ứng dụng, hệ thống đã thực hiện được các luồng chính của một tài liệu, bao gồm tạo tài liệu, cập nhật metadata, upload phiên bản file, gửi duyệt, phê duyệt, từ chối, xuất bản, preview/download, retention và lưu trữ. Mỗi tài liệu được quản lý cùng metadata, version, checksum, trạng thái workflow, lịch sử phê duyệt và audit event liên quan. Các audit log kết hợp với workflow history, retention evidence và evidence packet giúp kiểm tra lại tài liệu đã được xử lý như thế nào, bởi ai và theo kết quả nào.

Về bảo mật ứng dụng, nhóm đã tích hợp các cơ chế như RBAC, ACL, classification, Compliance Officer metadata-only boundary, grant token ngắn hạn, malware scanning, DLP và audit hash-chain. Những cơ chế này giúp kiểm soát ai được thao tác trên tài liệu, thao tác nào được phép, nội dung file được bảo vệ như thế nào và audit log có thể được kiểm tra tính toàn vẹn ra sao.

Ở phần DevSecOps, hệ thống đã được triển khai qua pipeline Jenkins với các bước kiểm thử, quét bảo mật, build image, đẩy artifact lên Harbor, cập nhật GitOps, đồng bộ qua Argo CD và kiểm tra sau triển khai. Các bằng chứng như scan report, image tag/digest, SBOM, trạng thái ký, GitOps revision, Argo CD health check, smoke test, ZAP report và dashboard quan sát vận hành giúp nhóm đối chiếu lại quá trình phát hành từ mã nguồn đến workload.

Nhìn chung, hệ thống đã đáp ứng được mục tiêu chính của đề tài ở mức đồ án chuyên ngành. DocVault chưa phải là một hệ thống production hoàn chỉnh, nhưng đã thể hiện được cách kết hợp giữa một ứng dụng quản lý tài liệu bảo mật và một quy trình DevSecOps có kiểm soát.

\section{Mức độ đáp ứng mục tiêu}

Đối chiếu với các mục tiêu ban đầu, đề tài đã đạt được các kết quả sau:

\begin{itemize}
  \item Xây dựng kiến trúc microservices với ranh giới trách nhiệm tương đối rõ giữa các dịch vụ.
  \item Hoàn thiện các luồng nghiệp vụ chính của hệ thống quản lý tài liệu và áp dụng kiểm soát truy cập ở backend.
  \item Cung cấp bằng chứng kiểm thử cho xác thực, RBAC/ACL, truy cập tài liệu nhạy cảm, malware, DLP, retention, audit và ranh giới quyền của Compliance Officer.
  \item Triển khai hạ tầng cloud-native trên EKS bằng Terraform, Helm, Kustomize và GitOps.
  \item Tích hợp các cổng chất lượng và bảo mật vào pipeline, bao gồm kiểm tra mã nguồn, dependency, secret, IaC, policy Kubernetes và container artifact.
  \item Áp dụng S3/KMS, IRSA, External Secrets, SBOM, ký artifact, backup và observability trong môi trường triển khai.
\end{itemize}

\section{Hạn chế}

Do thời gian và tài nguyên triển khai còn giới hạn, DocVault hiện tập trung vào việc chứng minh các luồng chính của hệ thống và quy trình DevSecOps. Phần ứng dụng đã thể hiện được các chức năng quan trọng như quản lý vòng đời tài liệu, phân quyền, audit, evidence và retention. Tuy nhiên, một số nội dung vẫn cần hoàn thiện nếu muốn triển khai trong môi trường production.

Ở phần ứng dụng, hệ thống vẫn cần bổ sung thêm kiểm thử trình duyệt đầy đủ, tối ưu trải nghiệm người dùng, hoàn thiện thêm các trường hợp biên của workflow và mở rộng khả năng xử lý khi số lượng tài liệu lớn hơn. Các cơ chế như DLP, malware scanning và retention hiện đã được tích hợp ở mức phục vụ đồ án, nhưng vẫn cần được đánh giá sâu hơn nếu áp dụng trong môi trường thực tế.

Ở phần DevSecOps, một số nội dung production hardening như NetworkPolicy chi tiết, admission controller bắt buộc xác minh chữ ký image, provenance chuẩn hóa, log persistence, backup/restore và diễn tập khôi phục dữ liệu vẫn chưa được hoàn thiện đầy đủ. Pipeline hiện đã chứng minh được luồng kiểm thử, scan, build, GitOps và kiểm tra sau triển khai, nhưng vẫn cần tối ưu thêm thời gian chạy và cách xử lý kết quả scan nếu áp dụng lâu dài.

\section{Hướng phát triển}

Trong thời gian tới, hệ thống có thể được phát triển tiếp theo hai hướng chính: hoàn thiện phần ứng dụng quản lý tài liệu và gia cố phần DevSecOps/hạ tầng triển khai.

Đối với phần ứng dụng, nhóm có thể mở rộng thêm các chức năng phục vụ quản lý tài liệu trong thực tế như tìm kiếm nâng cao, dashboard theo trạng thái tài liệu, báo cáo hoạt động theo người dùng, nhắc nhở tài liệu chờ duyệt, cải thiện evidence packet và bổ sung thêm các chính sách retention linh hoạt hơn. Ngoài ra, hệ thống cũng có thể được tối ưu thêm về giao diện, hiệu năng truy vấn và khả năng xử lý khi số lượng tài liệu, audit event và version file tăng lên.

Đối với phần bảo mật ứng dụng, nhóm có thể tiếp tục hoàn thiện DLP, malware scanning, kiểm soát tải xuống theo ngữ cảnh và kiểm thử các trường hợp biên của policy. Phần audit có thể mở rộng theo hướng lưu trữ bền vững hơn, tích hợp hệ thống log chuyên dụng hoặc bổ sung cơ chế ký số cho các evidence quan trọng.

Đối với phần DevSecOps, hướng phát triển tiếp theo là hoàn thiện các kiểm soát gần với production hơn như NetworkPolicy chi tiết, admission policy bắt buộc, xác minh chữ ký image, provenance chuẩn hóa, backup/restore, diễn tập khôi phục dữ liệu và tối ưu pipeline. Bên cạnh đó, hệ thống có thể bổ sung thêm cảnh báo vận hành, dashboard chi tiết hơn và kiểm thử tự động ở mức trình duyệt để tăng độ tin cậy trước khi phát hành.

\section{Kết luận chung}

Qua quá trình thực hiện đề tài, nhóm đã xây dựng và triển khai được hệ thống DocVault theo hai hướng chính: xây dựng ứng dụng quản lý tài liệu bảo mật và xây dựng quy trình DevSecOps cho hệ thống. Phần ứng dụng giúp quản lý tài liệu theo vòng đời, kiểm soát truy cập và ghi nhận audit/evidence. Phần DevSecOps giúp kiểm thử, quét bảo mật, đóng gói, triển khai và kiểm tra hệ thống theo một quy trình có thể lặp lại.

Ở phần ứng dụng, DocVault đã thể hiện được cách quản lý một tài liệu từ lúc tạo mới, upload version, gửi duyệt, phê duyệt, xuất bản cho đến retention hoặc archive. Tài liệu trong hệ thống không chỉ là một file lưu trữ, mà còn có metadata, phiên bản, trạng thái workflow, quyền truy cập, lịch sử phê duyệt, audit log và evidence. Nhờ đó, hệ thống có thể kiểm tra lại quá trình xử lý tài liệu khi cần.

Ở phần DevSecOps, nhóm đã xây dựng được pipeline để chạy kiểm thử, quét bảo mật, build image, đẩy artifact lên Harbor, cập nhật GitOps, triển khai qua Argo CD và kiểm tra sau triển khai. Các bằng chứng như log pipeline, report scan, image digest, SBOM, GitOps revision, trạng thái Argo CD, smoke test và ZAP report giúp đối chiếu lại quá trình phát hành hệ thống.

Trong quá trình thực hiện, nhóm cũng gặp nhiều khó khăn thực tế. Khi chạy local, các service có thể hoạt động tương đối đơn giản, nhưng khi đưa lên Kubernetes thì phát sinh nhiều vấn đề về biến môi trường, secret, kết nối nội bộ giữa các service và trạng thái readiness của workload. Việc cấu hình External Secrets, IRSA, Harbor, Argo CD và các manifest triển khai cũng mất nhiều thời gian để kiểm tra và đồng bộ.

Pipeline DevSecOps cũng phát sinh nhiều vấn đề trong quá trình chạy thử. Một số bước quét bảo mật làm pipeline chạy lâu hơn dự kiến, có lúc công cụ scan bị treo hoặc trả về cảnh báo CVE cần xem xét lại image và dependency. Việc xử lý các kết quả scan cũng giúp nhóm hiểu rõ hơn rằng DevSecOps không chỉ là thêm nhiều công cụ vào pipeline, mà còn cần biết kết quả nào quan trọng, kết quả nào có thể chấp nhận trong phạm vi demo và kết quả nào cần xử lý trước khi triển khai.

Thông qua đề tài, nhóm học được cách thiết kế một hệ thống microservices có nhiều thành phần phối hợp với nhau, cách tách trách nhiệm giữa metadata-service, document-service, workflow-service và audit-service, cũng như cách triển khai hệ thống lên môi trường Kubernetes. Nhóm cũng hiểu rõ hơn vai trò của audit log, evidence, image digest, SBOM, GitOps revision và các bằng chứng vận hành trong việc kiểm tra lại một hệ thống sau khi triển khai.

Tuy nhiên, hệ thống vẫn còn một số hạn chế. Một số phần như NetworkPolicy chi tiết, admission controller bắt buộc xác minh chữ ký image, provenance chuẩn hóa, kiểm thử trình duyệt đầy đủ, log persistence và diễn tập khôi phục dữ liệu vẫn cần tiếp tục hoàn thiện nếu muốn tiến gần hơn đến môi trường production. Ngoài ra, một số chức năng trong hệ thống hiện vẫn tập trung vào việc chứng minh luồng xử lý chính, chưa tối ưu sâu về hiệu năng, trải nghiệm người dùng và khả năng mở rộng trong môi trường thực tế.

Nhìn chung, kết quả hiện tại đã đáp ứng được mục tiêu của đồ án chuyên ngành. DocVault đã kết hợp được hai phần quan trọng: một ứng dụng quản lý tài liệu bảo mật có vòng đời, audit và evidence; cùng với một quy trình DevSecOps giúp kiểm thử, triển khai và kiểm chứng hệ thống. Đây là nền tảng để nhóm có thể tiếp tục mở rộng hệ thống theo hướng hoàn thiện hơn trong tương lai.
